\makeatletter
\def\input@path{{./}{../}{../../}{../../../}{../../../../}{preamble/}{../preamble/}{../../preamble/}{../../../preamble/}}
\makeatother

\input{preamble_exams.tex}
\usepackage{needspace}
\usepackage{titling}
\setlength{\droptitle}{-2em}
\pretitle{\par\begingroup\centering\LARGE}
\posttitle{\par\endgroup}
\preauthor{}\postauthor{}
\predate{}\postdate{}
\maketitlehookd{\vspace{-1em}}
\linespread{1.12}
\setlength{\parskip}{5pt}
\newgeometry{margin=0.65in,top=0.55in,bottom=0.65in}

\title{Grade 11 Term 2 Home Practice -- C3\\\large Anatomy of a confidence interval\\[0.25em]
	\normalsize Global Economics -- 2026--2027}
\author{}
\date{}

\begin{document}
	\maketitle

	\setcounter{tocdepth}{1}
	\setcounter{secnumdepth}{2}
	\phantomsection
	\label{toc}
	\tableofcontents

	\section*{Evaluated criteria}
	\begin{enumerate}[label=\textbf{C\arabic*},start=3]
		\item Identifies the components of a confidence interval, including point estimate, critical value, standard error, and margin of error, explaining how each component affects precision and reliability.
	\end{enumerate}

	\textbf{How to practise.} Try each task before opening its blue solution link. The arrows return to the contents, problem list or current problem.

	\section{Worked Solution Examples}
	\label{sec:worked-solution-examples}

	\SectionProblemsTableOfContents{
		\SectionProblemLink{prob:c3-1}{Household grocery spending (C3)} \\
		\SectionProblemLink{prob:c3-2}{Two estimates of average freight cost (C3)}
	}

	% ----------------------------------------------------------------
	\subsection{Household grocery spending (C3)}
	\label{prob:c3-1}

	An analyst estimates the mean weekly grocery bill for households in a district. The random sample mean is 84 thousand COP and the standard error of the sample mean is 3 thousand COP. A two-sided 95\% normal interval uses $z^*=1.96$. A two-sided 90\% interval uses $z^*=1.645$.

	\QuestionsTasks[sec:worked-solution-examples]{
		{3}/{sol:c3-1-a}/
		{Write the interval in the form estimate $\pm$ critical value $\times$ standard error. Identify all four components, including the margin of error, and explain what each contributes. Sketch the standard normal curve and mark the central area associated with the critical value.},
		{3}/{sol:c3-1-b}/
		{Recalculate the margin of error for 90\% confidence using the same data. Explain what happens to the centre, width and long-run coverage.}
	}

	\clearpage

	\Needspace{5\baselineskip}
	\phantomsection
	\label{sol:c3-1-a}
	\subsubsection*{C3 -- Identify the components and their roles}

	\noindent
	\begin{minipage}[t]{0.52\textwidth}
	\vspace{0pt}
	\centering
	\begin{tikzpicture}
	\begin{axis}[width=\linewidth,height=4.5cm,axis lines=left,
	xmin=-3.6,xmax=3.6,ymin=0,ymax=0.45,
	xlabel={$z$},ylabel={$\phi(z)$},xtick={-1.96,0,1.96},
	xticklabels={$-1.96$,$0$,$1.96$},ytick=\empty]
	\addplot[draw=none,fill=blue!30,domain=-1.96:1.96,samples=85]
	{exp(-x^2/2)/sqrt(2*pi)} \closedcycle;
	\addplot[black,thick,domain=-3.6:3.6,samples=125] {exp(-x^2/2)/sqrt(2*pi)};
	\addplot[black,dashed] coordinates {(-1.96,0) (-1.96,0.05844)};
	\addplot[black,dashed] coordinates {(1.96,0) (1.96,0.05844)};
	\node at (axis cs:0,0.18) {$0.95$};
	\end{axis}
	\end{tikzpicture}
	\end{minipage}
	\hfill
	\begin{minipage}[t]{0.44\textwidth}
	\vspace{0pt}
	The \textbf{point estimate} 84 sets the centre. The \textbf{standard error} 3 measures how much the sample mean typically varies across samples. The \textbf{critical value} 1.96 is a unitless multiplier. On the standard normal curve, the central area is about 0.95, leaving 0.025 in each tail.

	The \textbf{margin of error} is $E=1.96(3)=5.88$, the distance from the centre to each endpoint. The interval is $84\pm 5.88$, or $[78.12,89.88]$ thousand COP.
	\end{minipage}

	\secInicio[prob:c3-1][sec:worked-solution-examples]

	\Needspace{5\baselineskip}
	\phantomsection
	\label{sol:c3-1-b}
	\subsubsection*{C3 -- Separate confidence from precision}

	With the same sample, the centre stays at 84 and $SE$ stays at 3. The smaller critical value gives $E=1.645(3)=4.935$ thousand COP. The width drops from $2(5.88)=11.76$ to $2(4.935)=9.87$ thousand COP. The interval is more precise in the sense of being narrower, and its nominal long-run coverage falls from 95\% to 90\%.

	\secInicio[prob:c3-1][sec:worked-solution-examples]

	\clearpage

	% ----------------------------------------------------------------
	\subsection{Two estimates of average freight cost (C3)}
	\label{prob:c3-2}

	Two simple random samples estimate the mean freight cost per container on a route. Both sample means are 250 USD. Sample A gives $SE=8$ USD and sample B gives $SE=4$ USD. Both use the two-sided critical value $z^*=1.96$ for 95\% confidence.

	\QuestionsTasks[sec:worked-solution-examples]{
		{3}/{sol:c3-2-a}/
		{Find both margins and widths. Identify which components stay fixed, and explain whether the smaller standard error changes nominal coverage.},
		{3}/{sol:c3-2-b}/
		{Suppose sample B's point estimate had been 254 USD and its SE still 4 USD. Give the new interval and describe how it changes.}
	}

	\Needspace{5\baselineskip}
	\phantomsection
	\label{sol:c3-2-a}
	\subsubsection*{C3 -- Explain the effect of a smaller standard error}

	The point estimate stays at 250 USD and the critical value stays at 1.96. For A, $E=1.96(8)=15.68$ USD, giving $[234.32,265.68]$ and width 31.36 USD. For B, $E=1.96(4)=7.84$ USD, giving $[242.16,257.84]$ and width 15.68 USD. With the same critical value, halving the standard error halves the margin and width. Both procedures target 95\% coverage under their sampling and model assumptions.

	\secInicio[prob:c3-2][sec:worked-solution-examples]

	\Needspace{5\baselineskip}
	\phantomsection
	\label{sol:c3-2-b}
	\subsubsection*{C3 -- Explain the effect of moving the point estimate}

	The new interval is $254\pm 7.84=[246.16,261.84]$ USD. Both endpoints move 4 USD because the centre moves 4 USD. The same critical value and SE give the same margin and width, so the shift changes the estimated location while precision and nominal coverage stay the same.

	\secInicio[prob:c3-2][sec:worked-solution-examples]


\end{document}
